%\clearpage
\section{Continuidad uniforme}

\begin{remark}
Sea $f:X\to Y$ una función continua en $A\subseteq X$.
Si $p\in A$ y $\varepsilon>0$, existe $\delta_{p,\varepsilon}>0$ (que depende de $p$ y de $\varepsilon$)
tal que
\[
d_Y(f(p),f(x))<\varepsilon \quad \text{si } x\in B_X(p,\delta_{p,\varepsilon})\cap A.
\]
\end{remark}

\begin{definicion}[Continuidad uniforme]
Sea $f:X\to Y$. Decimos que $f$ es uniformemente continua en $A\subseteq X$ si para todo $\varepsilon>0$ existe $\delta_\varepsilon>0$ (depende sólo de $\varepsilon$) tal que para todo $x,y\in A$, si
\[
d_X(x,y)<\delta_\varepsilon,
\]
entonces
\[
d_Y(f(x),f(y))<\varepsilon.
\]
\end{definicion}

\begin{remark}
Si $f$ es uniformemente continua en $A$, entonces $f$ es continua en $A$.
\end{remark}

\begin{proof}
Sea $\varepsilon>0$. Como $f$ es uniformemente continua en $A$, existe $\delta>0$ tal que para todo $x,y\in A$, si $d_X(x,y)<\delta$, entonces $d_Y(f(x),f(y))<\varepsilon$.

Fijado $p\in A$, tomando $y=p$, se obtiene que para todo $x\in A$, si $d_X(x,p)<\delta$, entonces
\[
d_Y(f(x),f(p))<\varepsilon.
\]
Luego $f$ es continua en $p$, y como $p$ fue arbitrario, $f$ es continua en $A$.
\end{proof}

\begin{example}
La función $f(x)=x$ es uniformemente continua en $\R$.

En efecto, dado $\varepsilon>0$, tomando $\delta=\varepsilon$, si $|x-y|<\delta$ entonces
\[
|f(x)-f(y)|=|x-y|<\varepsilon.
\]
\end{example}

\begin{example}
La función $f(x)=x^2$ no es uniformemente continua en $\R$.

Supongamos que sí lo es. Con $\varepsilon=1$, existe $\delta>0$ tal que para todo $x,y\in\R$, si $|x-y|<\delta$ entonces $|f(x)-f(y)|<1$.

Escogemos $n\in\N$ tal que
\[
n\delta+\frac{\delta^2}{4}>1.
\]
Tomamos $x=n+\frac{\delta}{2}$ y $y=n$. Entonces
\[
|x-y|=\frac{\delta}{2}<\delta,
\]
pero
\[
|f(x)-f(y)|=\left|\left(n+\frac{\delta}{2}\right)^2-n^2\right|=n\delta+\frac{\delta^2}{4}>1,
\]
contradicción.
\end{example}

\begin{teorema}[Teorema de Heine]
Sea $f:X\to Y$ una función. Si $A\subseteq X$ es compacto y $f$ es continua en $A$, entonces $f$ es uniformemente continua en $A$.
\end{teorema}

\begin{proof}
Sea $\varepsilon>0$. Como $f$ es continua en $A$, para cada $p\in A$ existe $\delta_p>0$ tal que
\[
x\in B_X(p,\delta_p)\cap A \implies d_Y(f(p),f(x))<\frac{\varepsilon}{2}.
\]

La familia
\[
\mathcal{C}=\left\{B_X\left(p,\frac{\delta_p}{2}\right):p\in A\right\}
\]
es un cubrimiento abierto de $A$. Por compacidad, existen $p_1,\dots,p_k\in A$ tales que
\[
A\subseteq\bigcup_{i=1}^k B_X\left(p_i,\frac{\delta_{p_i}}{2}\right).
\]

Definimos
\[
\delta=\min\left\{\frac{\delta_{p_i}}{2}:i=1,\dots,k\right\}>0.
\]
Sean $x,y\in A$ con $d_X(x,y)<\delta$. Existe $j$ tal que
\[
x\in B_X\left(p_j,\frac{\delta_{p_j}}{2}\right).
\]
Entonces, por desigualdad triangular,
\[
d_X(y,p_j)\le d_X(y,x)+d_X(x,p_j)<\delta+\frac{\delta_{p_j}}{2}\le\delta_{p_j},
\]
así que $y\in B_X(p_j,\delta_{p_j})$.

Por tanto,
\[
d_Y(f(x),f(y))\le d_Y(f(x),f(p_j))+d_Y(f(y),f(p_j))<\frac{\varepsilon}{2}+\frac{\varepsilon}{2}=\varepsilon.
\]
Concluimos que $f$ es uniformemente continua en $A$.
\end{proof}

\begin{definicion}
Sea $f:X\to X$ una función.
\begin{enumerate}
    \item Un punto $p\in X$ es un punto fijo de $f$ si $f(p)=p$.
    \item Decimos que $f$ es una contracción de $X$ si existe $0<\alpha<1$ tal que
    \[
    d(f(x),f(y))\le \alpha d(x,y),\quad \forall x,y\in X.
    \]
\end{enumerate}
\end{definicion}

\begin{remark}
Si $f:X\to X$ es una contracción, entonces $f$ es uniformemente continua en $X$.
\end{remark}

\begin{teorema}[Teorema del punto fijo]
Sea $f:X\to X$ una contracción en un espacio métrico completo $X$.
Entonces existe un único punto $p\in X$ tal que $f(p)=p$.
\end{teorema}
